\scanpage{38}
\mspage{18}
que si on choisit un repère direct $r_0$ \deleted{\unsure{$\in SR$}} de $\deleted{\unsure{B}}SP_M$ -- ce qui revient à se donner une base symplectique $M\simeq\mathbb Z^2$ --, alors pour toute structure 2-polyédrale triangulée orientée $P$ et \unsure{tout} repère $r$ de $P$, il existe un homomorphisme et un seul
\[
\left\{\begin{aligned}
SF_M&\longrightarrow P&&\text{tel que}\\
r_0&\longmapsto r&&.
\end{aligned}\right.
\]

\(\overline{\underline{\mathrm{III}}}\)

Quand on se donne un module libre \inserted{$M$} de rang 2 sur un anneau $\Lambda$ \inserted{\deleted{\unsure{``ab}} ``absolu'' i.e.} isomorphe à \deleted{\unsure{un}} $\mathbb Z$ ou à un $\mathbb Z/n\mathbb Z$ (où on suppose néanmoins $n\geqslant3$), sans se donner de structure symplectique \inserted{$e_M$} ou semi-symplectique $\{\pm e_M\}$\,; on peut néanmoins lui associer encore un 2-polyèdre \inserted{triangulé} régulier, \deleted{de la} \inserted{dont \unsure{le|l'} $\pi_0$ est} paramétré par l'ensemble quotient
\[
(\mathrm{dét}\,M)^*/\pm1\simeq\mathrm{Ssympl}(M).\tag*{(89)}
\]
formé des structures $\tfrac12$ symplectiques sur $M$,
\margin{On a donc changé la notation p.r. à (23), qui finalement n'introduit pas ``le bon'' objet\ldots}
% slanted
en prenant le \inserted{polyèdre $P_M$} somme directe des 2-polyèdres réguliers triangulés (\uline{non} \uline{orientés}) associés aux diverses structures semisymplectiques sur $M$.

Notons que pour
\[
\mu\in(\mathrm{det}\,M)^*/\{\pm1\}\tag*{(90)}
\]
% archive sheet number: 37
